\section{AHSME 1995}
        \begin{enumerate}
        	\item (\textit{AHSME 1995}) Kim earned scores of 87, 83, and 88 on her first three mathematics examinations. If Kim receives a score of 90 on the fourth exam, then her average will 


 $\mathrm{(A) \ \text{remain the same} } \qquad \mathrm{(B) \ \text{increase by 1} } \qquad \mathrm{(C) \ \text{increase by 2} } \qquad \mathrm{(D) \ \text{increase by 3} } \qquad \mathrm{(E) \ \text{increase by 4} }$

 \textbf{Answer: }B



\textbf{Solution}

The average of the first three test scores is $\frac{88+83+87}{3}=86$. The average of all four exams is $\frac{87+83+88+90}{4}=87$. It increased by one point. $\mathrm{(B)}$



	\item (\textit{AHSME 1995}) If $\sqrt {2 + \sqrt {x}} = 3$, then $x =$
$\mathrm{(A) \ 1 } \qquad \mathrm{(B) \ \sqrt{7} } \qquad \mathrm{(C) \ 7 } \qquad \mathrm{(D) \ 49 } \qquad \mathrm{(E) \ 121 }$

 \textbf{Answer: }D



\textbf{Solution}

$2 + \sqrt{x} = 9 \Longrightarrow \sqrt{x} = 7 \Longrightarrow x = 49 \Rightarrow \mathrm{(D)}$



	\item (\textit{AHSME 1995}) The total in-store price for an appliance is $\textdollar 99.99$. A television commercial advertises the same product for three easy payments of $\textdollar 29.98$ and a one-time shipping and handling charge of $\textdollar 9.98$. How many cents are saved by buying the appliance from the television advertiser?


 $\mathrm{(A) \ 6 } \qquad \mathrm{(B) \ 7 } \qquad \mathrm{(C) \ 8 } \qquad \mathrm{(D) \ 9 } \qquad \mathrm{(E) \ 10 }$

 \textbf{Answer: }B



\textbf{Solution}

IMPORTANT NOTICE: The original problem statement had "how much is saved". However, because this made little sense when the calculations were done, the problem statement was changed to "how many cents".

 We see that 3 payments of $\textdollar 29.98$ will be a total cost of $3\cdot(30-.02)=90-.06$

 Adding this to $\textdollar9.98$ we have a total of $99.98-.06$

 Clearly, this differs from $\textdollar 99.99$ by $7$ cents. Thus, the answer is $\fbox{\text{(B)}}$



	\item (\textit{AHSME 1995}) If $M$ is $30 \%$ of $Q$, $Q$ is $20 \%$ of $P$, and $N$ is $50 \%$ of $P$, then $\frac {M}{N} =$
$\mathrm{(A) \ \frac {3}{250} } \qquad \mathrm{(B) \ \frac {3}{25} } \qquad \mathrm{(C) \ 1 } \qquad \mathrm{(D) \ \frac {6}{5} } \qquad \mathrm{(E) \ \frac {4}{3} }$

 \textbf{Answer: }B



\textbf{Solution 1}

We are given: $M=\frac{3Q}{10}$, $Q=\frac{P}{5}$, $N=\frac{P}{2}$. We want M in terms of N, so we substitute N into everything:

 $\frac{2}{5}N=\frac{P}{5}=Q$
$M=\frac{3N}{25}$
$\frac{M}{N}=\frac{3}{25} \Rightarrow \mathrm{(B)}$




\textbf{Solution 2}

Alternatively, picking an arbitrary value for $Q$ of $100$, we find that $M = 30\% \cdot 100 = 0.30 \cdot 100 = 30$.

 We find that $Q  = 20\% \cdot P$, meaning $100 = 0.2\cdot P$, giving $P = \frac{100}{0.2} = \frac{1000}{2} = 500$.

 Finally, since $N$ is $50\%$ of $P$, we have $N = 50\% \cdot 500 = 0.5 \cdot 500 = 250$.

 Thus, $M = 30$ and $N = 250$, so their ratio $\frac{M}{N} = \frac{30}{250} = \frac{3}{25} \Rightarrow \mathrm{(B)}$

 This method does not prove that the answer must be constant, but it proves that if the answer is a constant, it must be $B$.

 


	\item (\textit{AHSME 1995}) A rectangular field is 300 feet wide and 400 feet long. Random sampling indicates that there are, on the average, three ants per square inch through out the field. [12 inches = 1 foot.] Of the following, the number that most closely approximates the number of ants in the field is 


 $\mathrm{(A) \ \text{500 thousand} } \qquad \mathrm{(B) \ \text{5 million} } \qquad \mathrm{(C) \ \text{50 million} } \qquad \mathrm{(D) \ \text{500 million} } \qquad \mathrm{(E) \ \text{5 billion} }$

 \textbf{Answer: }C



\textbf{Solution}

The rectangular field is $300 \text{ feet} \cdot \frac{12 \text{ inches}}{1 \text{ foot}} =  3600 \text{ inches}$ wide and $400 \text{ feet} \cdot \frac{12 \text{ inches}}{1 \text{ foot}} =  4800 \text{ inches}$ inches long.

 It has an area of $3.6\cdot 10^3 \cdot 4.8 \cdot 10^3 = 17.28 \cdot 10^6$ square inches. 

 We multiply by $3$ to account for the ants to get approximately $50 \cdot 10^6$, which is $50$ million $\Rightarrow \mathrm{(C)}$



	\item (\textit{AHSME 1995}) The figure shown can be folded into the shape of a cube. In the resulting cube, which of the lettered faces is opposite the face marked $x$? 


 
                    \begin{figure}[H]
                        \centering
                        \includegraphics[width=0.5\textwidth]{Problem/AHSME_I_1995_6_Problem_Figure_1.png}
                    \end{figure}
                    
$\mathrm{(A) \ A } \qquad \mathrm{(B) \ B } \qquad \mathrm{(C) \ C } \qquad \mathrm{(D) \ D } \qquad \mathrm{(E) \ E }$

 \textbf{Answer: }C



\textbf{Solution 1}

The marked side is the side with the $x.$ We imagine it folding up. First, we fold the $x$ upwards. Now we fold the $A$ upwards, and thus $x$ is touching $B$ on it's left side. We now fold $B$ up, and we realize that $x$ won't be touching $\boxed{\mathrm{(C)}}$ at all.

 



\textbf{Solution 2 (Cheese)}

Cut out this shape on your scratch paper. Label the letters and $x$. Fold it and get the answer $\boxed{\mathrm{(C)}}$. 

 \textit{Note:} If cutting paper is not allowed, just ask for $6$ pieces of paper and do it. It's gonna be harder though.

 


	\item (\textit{AHSME 1995}) The radius of Earth at the equator is approximately 4000 miles. Suppose a jet flies once around Earth at a speed of 500 miles per hour relative to Earth. If the flight path is a negligible height above the equator, then, among the following choices, the best estimate of the number of hours of flight is: 


 $\mathrm{(A) \ 8 } \qquad \mathrm{(B) \ 25 } \qquad \mathrm{(C) \ 50 } \qquad \mathrm{(D) \ 75 } \qquad \mathrm{(E) \ 100 }$

 \textbf{Answer: }C



\textbf{Solution}

We want the number of hours that it takes the jet to fly the length of the circumference. $\frac{8000\pi}{500}=16\pi$. The best estimate of that is $50\Rightarrow \mathrm{(C)}$



	\item (\textit{AHSME 1995}) In $\triangle ABC$, $\angle C = 90^\circ, AC = 6$ and $BC = 8$. Points $D$ and $E$ are on $\overline{AB}$ and $\overline{BC}$, respectively, and $\angle BED = 90^\circ$. If $DE = 4$, then $BD =$

                    \begin{figure}[H]
                        \centering
                        \includegraphics[width=0.5\textwidth]{Problem/AHSME_I_1995_8_Problem_Figure_1.png}
                    \end{figure}
                    
$\mathrm{(A) \ 5 } \qquad \mathrm{(B) \ \frac {16}{3} } \qquad \mathrm{(C) \ \frac {20}{3} } \qquad \mathrm{(D) \ \frac {15}{2} } \qquad \mathrm{(E) \ 8 }$

 \textbf{Answer: }C



\textbf{Solution}


                    \begin{figure}[H]
                        \centering
                        \includegraphics[width=0.5\textwidth]{Solution/AHSME_I_1995_8_Solution_Figure_1.png}
                    \end{figure}
                    
$\triangle BAC$ is a $6-8-10$ right triangle with hypotenuse $AB = 10$.

 $\triangle BDE$ is similar to $\triangle BAC$ by angle-angle similarity since $E=C = 90^\circ$ and $B=B$, and thus $\frac{BD}{BA} = \frac{DE}{AC}$.

 Solving the above for $BD$, we get $BD=\frac{BA\cdot DE}{AC} = 10\cdot \dfrac{4}{6}=\dfrac{20}{3}\Rightarrow \boxed{\mathrm{(C)}}$.

 


	\item (\textit{AHSME 1995}) Consider the figure consisting of a square, its diagonals, and the segments joining the midpoints of opposite sides. The total number of triangles of any size in the figure is 

                    \begin{figure}[H]
                        \centering
                        \includegraphics[width=0.5\textwidth]{Problem/AHSME_I_1995_9_Problem_Figure_1.png}
                    \end{figure}
                    
$\mathrm{(A) \ 10 } \qquad \mathrm{(B) \ 12 } \qquad \mathrm{(C) \ 14 } \qquad \mathrm{(D) \ 16 } \qquad \mathrm{(E) \ 18 }$

 \textbf{Answer: }D



\textbf{Solution}


                    \begin{figure}[H]
                        \centering
                        \includegraphics[width=0.5\textwidth]{Solution/AHSME_I_1995_9_Solution_Figure_1.png}
                    \end{figure}
                    

 There are 8 little triangles, 4 triangles with twice the area, and 4 triangles with four times the area of the smaller triangles. $8+4+4=16\Rightarrow \mathrm{(D)}$



	\item (\textit{AHSME 1995}) The area of the triangle bounded by the lines $y = x, y = - x$ and $y = 6$ is


 $\mathrm{(A) \ 12 } \qquad \mathrm{(B) \ 12\sqrt{2} } \qquad \mathrm{(C) \ 24 } \qquad \mathrm{(D) \ 24\sqrt{2} } \qquad \mathrm{(E) \ 36 }$

 \textbf{Answer: }E



\textbf{Solution}


                    \begin{figure}[H]
                        \centering
                        \includegraphics[width=0.5\textwidth]{Solution/AHSME_I_1995_10_Solution_Figure_1.png}
                    \end{figure}
                    

 The height of the triangle is $y = 6$, and the width of the triangle is $|x_1| + |x_2| = 2y = 12$. Thus the area of the triangle is $\frac 12 \cdot 6 \cdot 12 = 36\ \mathrm{(E)}$.

 


	\item (\textit{AHSME 1995}) How many base 10 four-digit numbers, $N = \underline{a} \underline{b} \underline{c} \underline{d}$, satisfy all three of the following conditions?


 $\text{(i)}$ $4,000 \leq N < 6,000;$
$\text{(ii)}$ $N$ $\text{is a multiple of 5}$; 


 $\text{(iii)}$ $3 \leq b < c \leq 6$.


 

$\mathrm{(A) \ 10 } \qquad \mathrm{(B) \ 18 } \qquad \mathrm{(C) \ 24 } \qquad \mathrm{(D) \ 36 } \qquad \mathrm{(E) \ 48 }$

 \textbf{Answer: }C



\textbf{Solution}


            \begin{itemize}
        \item For condition (i), the restriction is put on $a$; $N<4000$ if $a<4$, and $N \ge 6$ if $a \ge 6$. Therefore, $a=4,5$.
        \item For condition (ii), the restriction is put on $d$; it must be a multiple of $5$. Therefore, $d=0,5$.
        \item For condition (iii), the restriction is put on $b$ and $c$. The possible ordered pairs of $b$ and $c$ are $(3,4)$, $(3,5)$, $(3,6)$, $(4,5), (4,6),$ and $(5,6),$ and there are $6$ of them. Alternatively, we are picking from the four digits 3, 4, 5, 6, and for every combination of two, there is exactly one way to arrange them in increasing order, so we have $\binom{4}{2} = 6$ choices for $b$ and $c$ when we consider them together.
            \end{itemize}
            Multiplying the possibilities for each restriction, $2 \cdot 2 \cdot 6=24\Rightarrow \mathrm{(C)}$.

 


	\item (\textit{AHSME 1995}) Let $f$ be a linear function with the properties that $f(1) \leq f(2), f(3) \geq f(4),$ and $f(5) = 5$. Which of the following is true?


 $\mathrm{(A) \ f(0) < 0 } \qquad \mathrm{(B) \ f(0) = 0 } \qquad \mathrm{(C) \ f(1) < f(0) < f( - 1) } \qquad \mathrm{(D) \ f(0) = 5 } \qquad \mathrm{(E) \ f(0) > 5 }$

 \textbf{Answer: }D



\textbf{Solution 1}

A linear function has the property that $f(a)\leq f(b)$ either for all $a<b$, or for all $b<a$. Since $f(3)\geq f(4)$, $f(1)\geq f(2)$. Since $f(1)\leq f(2)$, $f(1)=f(2)$. And if $f(a)=f(b)$ for $a\neq b$, then $f(x)$ is a constant function. Since $f(5)=5$, $f(0)=5\Rightarrow \mathrm{(D)}$




\textbf{Solution 2}

If $f$ is a linear function, the statement $f(1) \leq f(2)$ states that the slope of the line $\frac{f(2) - f(1)}{2 - 1} = f(2) - f(1)$ is nonnegative:  it is either positive or zero.

 Similarly, the statement $f(3) \geq f(4)$ states that the slope of the line $\frac{f(4) - f(3)}{4 - 3} = f(4) - f(3)$ is nonpositive:  it is either negative or zero.

 Since the slope of a linear function can only have one value, it must be zero, and thus the function is a constant.  The statement $f(5) = 5$ tells us that the value of the constant is $5$, and thus that $f(x) = 5$.  This leads to $f(0)=5\Rightarrow \mathrm{(D)}$




\textbf{Solution 3 (Common Sense)}

It should be very clear that $f(1)<f(2)$ and $f(3)>f(4)$ is contradictory because of the fact that linear functions are monotonic. The only thing that makes sense is $f(1)=f(2)$, and $f(3)=f(4)$. This means that $f(x)$ has a slope of $0$. So $f(x)=5$. So $f(0)=5$. Select $\boxed{D}$. 

 


	\item (\textit{AHSME 1995}) The addition below is incorrect. The display can be made correct by changing one digit $d$, wherever it occurs, to another digit $e$. Find the sum of $d$ and $e$.


 $\begin{tabular}{ccccccc} & 7 & 4 & 2 & 5 & 8 & 6 \\ + & 8 & 2 & 9 & 4 & 3 & 0 \\ \hline 1 & 2 & 1 & 2 & 0 & 1 & 6 \end{tabular}$
$\mathrm{(A) \ 4 } \qquad \mathrm{(B) \ 6 } \qquad \mathrm{(C) \ 8 } \qquad \mathrm{(D) \ 10 } \qquad \mathrm{(E) \ \text{more than 10} }$

 \textbf{Answer: }C



\textbf{Solution}

If we change $0$, the units column would  be incorrect. 

 If we change $1$, then the leading $1$ in the sum would be incorrect.

 However, looking at the $2$ in the hundred-thousands column, it would be possible to change the $2$ to either a $5$ (no carry) or a $6$ (carry) to create a correct statement.  

 Changing the $2$ to a $5$ would give $745586 + 859430$ on top, which equals $1605016$.  This does not match up to the bottom.

 Changing the $2$ to a $6$ gives $746586 + 869430$ on top, which has a sum of $1616016$.  This is the number on the bottom if the $2$s were changed to $6$s.

 Thus $d=2$ and $e=6$. so $d+e= 8 \boxed{\mathrm{ (C)}}$



	\item (\textit{AHSME 1995}) If $f(x) = ax^4 - bx^2 + x + 5$ and $f( - 3) = 2$, then $f(3) =$
$\mathrm{(A) \ -5 } \qquad \mathrm{(B) \ -2 } \qquad \mathrm{(C) \ 1 } \qquad \mathrm{(D) \ 3 } \qquad \mathrm{(E) \ 8 }$

 \textbf{Answer: }E



\textbf{Solution 1}

$f(-x) = a(-x)^4 - b(-x)^2 - x + 5$
$f(-x) = ax^4 - bx^2 - x + 5$
$f(-x) = (ax^4 - bx^2 + x + 5) - 2x$
$f(-x) =  f(x) - 2x$. 

 Thus $f(3) = f(-3)-2(-3) = 8 \Rightarrow \mathrm{(E)}$.

 



\textbf{Solution 2}

If $f(-3) = 2$, then $81a - 9b + -3 + 5 = 2$.  Simplifying, we get $81a - 9b = 0$.

 Getting an expression for $f(3)$, we find $f(3) = 81a - 9b + 3 + 5$.  Since the first two terms sum up to zero, we get $f(3) = 8$, which is answer $\mathrm{(E)}$




\textbf{Solution 3}

Substituting $x = -3$, we get
 \[f(-3) = 81a - 9b - 3 + 5 = 81a - 9b + 2.\]But $f(-3) = 2$, so $81a - 9b + 2 = 2$, which means $81a - 9b = 0$. Then
 \[f(3) = 81a - 9b + 3 + 5 = 0 + 3 + 5 = \boxed{8}.\]



	\item (\textit{AHSME 1995}) Five points on a circle are numbered 1,2,3,4, and 5 in clockwise order. A bug jumps in a clockwise direction from one point to another around the circle; if it is on an odd-numbered point, it moves one point, and if it is on an even-numbered point, it moves two points. If the bug begins on point 5, after 1995 jumps it will be on point 

                    \begin{figure}[H]
                        \centering
                        \includegraphics[width=0.5\textwidth]{Problem/AHSME_I_1995_15_Problem_Figure_1.png}
                    \end{figure}
                    
$\mathrm{(A) \ 1 } \qquad \mathrm{(B) \ 2 } \qquad \mathrm{(C) \ 3 } \qquad \mathrm{(D) \ 4 } \qquad \mathrm{(E) \ 5 }$

 \textbf{Answer: }D



\textbf{Solution}

Let us see how the bug moves.

 First, we see that if it starts at point 5, it moves to point 1.

 At point 1, it moves to point 2.

 At point 2, since 2 is even, it moves to point 4.

 Then at point 4, it moves to point 1.

 We can see that at this point, the bug will cycle between 1, 2, and 4

 More specifically, we can see that all numbers congruent to 0 (mod 3) will have the bug on point 4 on that step number.

 Thus, we can conclude that the answer is $\fbox{\text{(D)}}$



	\item (\textit{AHSME 1995}) Anita attends a baseball game in Atlanta and estimates that there are 50,000 fans in attendance. Bob attends a baseball game in Boston and estimates that there are 60,000 fans in attendance. A league official who knows the actual numbers attending the two games notes that:


 i. The actual attendance in Atlanta is within $10 \%$ of Anita's estimate.
ii. Bob's estimate is within $10 \%$ of the actual attendance in Boston.


 To the nearest 1,000, the largest possible difference between the numbers attending the two games is


 $\mathrm{(A) \ 10000 } \qquad \mathrm{(B) \ 11000 } \qquad \mathrm{(C) \ 20000 } \qquad \mathrm{(D) \ 21000 } \qquad \mathrm{(E) \ 22000 }$

 \textbf{Answer: }E



\textbf{Solution}

Since the number of people at the game in Boston is certainly more than the number of fans in Atlanta, we need to compute the maximum of Bob's game minus the minimum of Anita's game. Note however that there is a slight different between conditions (i) and (ii); the attendence is \textit{within} $\pm 10 \%$ from Anita's estimate of $50,000$, but Bob's estimate is \textit{within} $\pm 10 \%$ of the actual attendance, or $0.9x \le 60000 \le 1.1x \Longrightarrow x \le \frac{60000}{0.9}$. 

 The answer is $\frac{60000}{0.9} - 50000 \times 0.9 \approx 22000 \Longrightarrow \mathrm{(E)}$.

 


	\item (\textit{AHSME 1995}) Given regular pentagon $ABCDE$, a circle can be drawn that is tangent to $\overline{DC}$ at $D$ and to $\overline{AB}$ at $A$. The number of degrees in minor arc $AD$ is


 
                    \begin{figure}[H]
                        \centering
                        \includegraphics[width=0.5\textwidth]{Problem/AHSME_I_1995_17_Problem_Figure_1.png}
                    \end{figure}
                    
$\mathrm{(A) \ 72 } \qquad \mathrm{(B) \ 108 } \qquad \mathrm{(C) \ 120 } \qquad \mathrm{(D) \ 135 } \qquad \mathrm{(E) \ 144 }$

 \textbf{Answer: }E



\textbf{Solution 1}

Define major arc DA as $DA$, and minor arc DA as $da$. Extending DC and AB to meet at F, we see that $\angle CFB=36=\frac{DA-da}{2}$. We now have two equations: $DA-da=72$, and $DA+da=360$. Solving, $DA=216$ and $da=144\Rightarrow \mathrm{(E)}$.

 


	\item (\textit{AHSME 1995}) Two rays with common endpoint $O$ form a $30^\circ$ angle. Point $A$ lies on one ray, point $B$ on the other ray, and $AB = 1$. The maximum possible length of $OB$ is


 $\mathrm{(A) \ 1 } \qquad \mathrm{(B) \ \frac {1 + \sqrt {3}}{\sqrt 2} } \qquad \mathrm{(C) \ \sqrt{3} } \qquad \mathrm{(D) \ 2 } \qquad \mathrm{(E) \ \frac{4}{\sqrt{3}} }$

 \textbf{Answer: }D



\textbf{Solution}

Triangle $OAB$ has the property that $\angle O=30^{\circ}$ and $AB=1$. 

 From the \href{https://artofproblemsolving.com/wiki/index.php?title=Law\_of\_Sines}{Law of Sines}, $\frac{\sin{\angle A}}{OB}=\frac{\sin{\angle O}}{AB}$. 

 Since $\sin 30^\circ = \frac{1}{2}$, we have:

 $\frac{\sin{\angle A}}{OB}=\frac{\frac{1}{2}}{1}$
$2\sin{\angle A}=OB$. 

 The maximum of $\sin{\angle A}$ is $1$ when $\angle A = 90^\circ$, so the maximum of $OB$ is $2\Rightarrow \mathrm{(D)}$.

 


	\item (\textit{AHSME 1995}) Equilateral triangle $DEF$ is inscribed in equilateral triangle $ABC$ such that $\overline{DE} \perp \overline{BC}$. The ratio of the area of $\triangle DEF$ to the area of $\triangle ABC$ is

                    \begin{figure}[H]
                        \centering
                        \includegraphics[width=0.5\textwidth]{Problem/AHSME_I_1995_19_Problem_Figure_1.png}
                    \end{figure}
                    
$\mathrm{(A) \ \frac {1}{6} } \qquad \mathrm{(B) \ \frac {1}{4} } \qquad \mathrm{(C) \ \frac {1}{3} } \qquad \mathrm{(D) \ \frac {2}{5} } \qquad \mathrm{(E) \ \frac {1}{2} }$

 \textbf{Answer: }C



\textbf{Solution}


                    \begin{figure}[H]
                        \centering
                        \includegraphics[width=0.5\textwidth]{Solution/AHSME_I_1995_19_Solution_Figure_1.png}
                    \end{figure}
                    

 Let's take one of the smaller right triangles. Without loss of generality, let the smaller leg be $1$. Since the triangle is a 30-60-90 right triangle, then the other leg is $\sqrt{3}$ and the hypotenuse is $2$. The side length of the bigger triangle is $1 + 2 = 3$ and the side length of the smaller triangle is $\sqrt{3}$. The ratio of the areas of two similar triangles is the square of the ratio of two corresponding side lengths, so the ratio of the area of triangle DEF to the area of triangle ABC is $\left(\frac{\sqrt{3}}{3}\right)^2 = \dfrac{1}{3} \rightarrow (C)$.

 


	\item (\textit{AHSME 1995}) If $a,b$ and $c$ are three (not necessarily different) numbers chosen randomly and with replacement from the set $\{1,2,3,4,5 \}$, the probability that $ab + c$ is even is


 $\mathrm{(A) \ \frac {2}{5} } \qquad \mathrm{(B) \ \frac {59}{125} } \qquad \mathrm{(C) \ \frac {1}{2} } \qquad \mathrm{(D) \ \frac {64}{125} } \qquad \mathrm{(E) \ \frac {3}{5} }$

 \textbf{Answer: }B



\textbf{Solution}

The probability of $ab$ being odd is $\left(\frac 35\right)^2 = \frac{9}{25}$, so the probability of $ab$ being even is $1 - \frac{9}{25} = \frac {16}{25}$.

 The probability of $c$ being odd is $3/5$ and being even is $2/5$
$ab+c$ is even if $ab$ and $c$ are both odd, with probability $\frac{9}{25} \cdot \frac{3}{5} = \frac{27}{125}$; or $ab$ and $c$ are both even, with probability $\frac{16}{25} \cdot \frac{2}{5} = \frac{32}{125}$. Thus the total probability is $\frac{59}{125} \Rightarrow \mathrm{(B)}$.

 


	\item (\textit{AHSME 1995}) Two nonadjacent vertices of a rectangle are $(4,3)$ and $(-4,-3)$, and the coordinates of the other two vertices are integers. The number of such rectangles is


 $\mathrm{(A) \ 1 } \qquad \mathrm{(B) \ 2 } \qquad \mathrm{(C) \ 3 } \qquad \mathrm{(D) \ 4 } \qquad \mathrm{(E) \ 5 }$

 \textbf{Answer: }E



\textbf{Solution}

The center of the rectangle is $(0,0)$, and the distance from the center to a corner is $\sqrt{4^2+3^2}=5$. The remaining two vertices of the rectangle must be another pair of points opposite each other on the circle of radius 5 centered at the origin. Let these points have the form $(\pm x,\pm y)$, where $x^2+y^2=25$. This equation has six pairs of integer solutions: $(\pm 4, \pm 3)$, $(\pm 4, \mp 3)$, $(\pm 3, \pm 4)$, $(\pm 3, \mp 4)$, $(\pm 5, 0)$, and $(0, \pm 5)$. The first pair of solutions are the endpoints of the given diagonal, and the other diagonal must span one of the other five pairs of points. $\Rightarrow \mathrm{(E)}$



	\item (\textit{AHSME 1995}) A pentagon is formed by cutting a triangular corner from a rectangular piece of paper. The five sides of the pentagon have lengths 13, 19, 20, 25 and 31, although this is not necessarily their order around the pentagon. The area of the pentagon is


 $\mathrm{(A) \ 459 } \qquad \mathrm{(B) \ 600 } \qquad \mathrm{(C) \ 680 } \qquad \mathrm{(D) \ 720 } \qquad \mathrm{(E) \ 745 }$

 \textbf{Answer: }E



\textbf{Solution}


                    \begin{figure}[H]
                        \centering
                        \includegraphics[width=0.5\textwidth]{Solution/AHSME_I_1995_22_Solution_Figure_1.png}
                    \end{figure}
                    

 Since the pentagon is cut from a rectangle, the cut-off triangle must be right. Since all of the lengths given are integers, it follows that this triangle is a \href{https://artofproblemsolving.com/wiki/index.php?title=Pythagorean\_Triple}{Pythagorean Triple}. We know that $31$ and either $25,\, 20$ must be the dimensions of the rectangle, since they are the largest lengths. With some trial-and-error, if we assign the shortest side, $13$, to be the \href{https://artofproblemsolving.com/wiki/index.php?title=Hypotenuse}{hypotenuse} of the triangle, we see the $5-12-13$ triple. Indeed this works, by placing the $31$ side opposite from the $19$ side and the $25$ side opposite from the $20$ side, leaving the cutaway side to be, as before, $13$.

 To find the area of the pentagon, we subtract the area of the triangle from that of the big rectangle: $31\cdot 25-\frac{12\cdot5}{2}=775-30=745\Longrightarrow \boxed{\mathrm{(E)}745}$.

 


	\item (\textit{AHSME 1995}) The sides of a triangle have lengths 11, 15, and $k$, where $k$ is an integer. For how many values of $k$ is the triangle obtuse?


 $\mathrm{(A) \ 5 } \qquad \mathrm{(B) \ 7 } \qquad \mathrm{(C) \ 12 } \qquad \mathrm{(D) \ 13 } \qquad \mathrm{(E) \ 14 }$

 \textbf{Answer: }D



\textbf{Solution}

By the \href{https://artofproblemsolving.com/wiki/index.php?title=Law\_of\_Cosines}{Law of Cosines}, a triangle is obtuse if the sum of the squares of two of the sides of the triangles is less than the square of the third. The largest angle is either opposite side $15$ or side $k$. If $15$ is the largest side, 

 
 \[15^2 >11^2 + k^2 \Longrightarrow k < \sqrt{104}\]

 By the \href{https://artofproblemsolving.com/wiki/index.php?title=Triangle\_Inequality}{Triangle Inequality} we also have that $k > 4$, so $k$ can be $5, 6, 7, \ldots , 10$, or $6$ values.

 If $k$ is the largest side, 

 
 \[k^2 >11^2 + 15^2 \Longrightarrow k > \sqrt{346}\]

 Combining with the Triangle Inequality $19 \le k < 26$, or $7$ values. These total $13\ \mathrm{(D)}$ values of $k$. 

 


	\item (\textit{AHSME 1995}) There exist positive integers $A,B$ and $C$, with no common factor greater than 1, such that


 
 \[A \log_{200} 5 + B \log_{200} 2 = C\]


 What is $A + B + C$?


 $\mathrm{(A) \ 6 } \qquad \mathrm{(B) \ 7 } \qquad \mathrm{(C) \ 8 } \qquad \mathrm{(D) \ 9 } \qquad \mathrm{(E) \ 10 }$

 \textbf{Answer: }A



\textbf{Solution}


 \[A \log_{200} 5 + B \log_{200} 2 = C\]

 Simplifying and taking the \href{https://artofproblemsolving.com/wiki/index.php?title=Logarithm}{logarithms} away,

 
 \[5^A \cdot 2^B=200^C=2^{3C} \cdot 5^{2C}\]

 Therefore, $A=2C$ and $B=3C$. Since $A, B,$ and $C$ are relatively prime, $C=1$, $B=3$, $A=2$. $A+B+C=6 \Rightarrow \mathrm{(A)}$



	\item (\textit{AHSME 1995}) A list of five positive integers has mean 12 and range 18. The mode and median are both 8. How many different values are possible for the second- largest element of the list? 


 $\mathrm{(A) \ 4 } \qquad \mathrm{(B) \ 6 } \qquad \mathrm{(C) \ 8 } \qquad \mathrm{(D) \ 10 } \qquad \mathrm{(E) \ 12 }$

 \textbf{Answer: }B



\textbf{Solution}

Let $a$ be the smallest element, so $a+18$ is the largest element. Since the mode is $8$, at least two of the five numbers must be $8$. The last number we denote as $b$.

 Then their average is $\frac{a + (8) + (8) + b + (a+18)}5 = 12 \Longrightarrow 2a + b = 26$. Clearly $a \le 8$. Also we have $b \le a + 18 \Longrightarrow 26-2a \le a + 18 \Longrightarrow 8/3 < 3 \le a$. Thus there are a maximum of $6$ values of $a$ which corresponds to $6$ values of $b$; listing shows that all such values work. The answer is $\boxed{B}$.

 


	\item (\textit{AHSME 1995}) In the figure, $\overline{AB}$ and $\overline{CD}$ are diameters of the circle with center $O$, $\overline{AB} \perp \overline{CD}$, and chord $\overline{DF}$ intersects $\overline{AB}$ at $E$. If $DE = 6$ and $EF = 2$, then the area of the circle is 


 
                    \begin{figure}[H]
                        \centering
                        \includegraphics[width=0.5\textwidth]{Problem/AHSME_I_1995_26_Problem_Figure_1.png}
                    \end{figure}
                    
$\mathrm{(A) \ 23 \pi } \qquad \mathrm{(B) \ \frac {47}{2} \pi } \qquad \mathrm{(C) \ 24 \pi } \qquad \mathrm{(D) \ \frac {49}{2} \pi } \qquad \mathrm{(E) \ 25 \pi }$

 \textbf{Answer: }C



\textbf{Solution}

\textbf{Solution 1}

 Let the radius of the circle be $r$ and let $x=\overline{OE}$. 

 By the \href{https://artofproblemsolving.com/wiki/index.php?title=Pythagorean\_Theorem}{Pythagorean Theorem}, $OD^2+OE^2=DE^2 \Rightarrow r^2+x^2=6^2=36$. 

 By \href{https://artofproblemsolving.com/wiki/index.php?title=Power\_of\_a\_point}{Power of a point}, $AE \cdot EB = DE \cdot EF \Rightarrow (r+x)(r-x)=r^2-x^2=6\cdot2=12$.  

 Adding these equations yields $2r^2=48 \Rightarrow r^2 = 24$. 

 Thus, the area of the circle is $\pi r^2 = 24\pi \Rightarrow C$.

 \textbf{Solution 2}

 Let the radius of the circle be $r$.

 We can see that $\triangle CFD$ has a right angle at $F$ and that $\triangle EOD$ has a right angle at $O$.  

 Both triangles also share $\angle ODE$, so $\triangle CFD$ and $\triangle EOD$ are \href{https://artofproblemsolving.com/wiki/index.php?title=Similar(geometry)&action=edit&redlink=1}{similar}.

 This means that $\frac {DE}{OD}=\frac {CD}{FD}$.

 So,  $\frac {6}{r}=\frac {2r}{8}$. Simplifying, $r^2 = 24$.

 This means the area is $24\pi \Rightarrow C$.

 


	\item (\textit{AHSME 1995}) Consider the triangular array of numbers with 0,1,2,3,... along the sides and interior numbers obtained by adding the two adjacent numbers in the previous row. Rows 1 through 6 are shown.


 
 \[\begin{tabular}{ccccccccccc} & & & & & 0 & & & & & \\ & & & & 1 & & 1 & & & & \\ & & & 2 & & 2 & & 2 & & & \\ & & 3 & & 4 & & 4 & & 3 & & \\ & 4 & & 7 & & 8 & & 7 & & 4 & \\ 5 & & 11 & & 15 & & 15 & & 11 & & 5  \end{tabular}\]


 Let $f(n)$ denote the sum of the numbers in row $n$. What is the remainder when $f(100)$ is divided by 100?


 $\mathrm{(A) \ 12 } \qquad \mathrm{(B) \ 30 } \qquad \mathrm{(C) \ 50 } \qquad \mathrm{(D) \ 62 } \qquad \mathrm{(E) \ 74 }$

 \textbf{Answer: }E



\textbf{Solution 1}

Note that if we re-draw the table with an additional diagonal row on each side, the table is actually just two of \href{https://artofproblemsolving.com/wiki/index.php?title=Pascal\%27s\_Triangle}{Pascal's Triangles}, except translated and summed. 
 \[\begin{tabular}{ccccccccccccccc} & & & & & 1 & & 0 & & 1 & & & & \\ & & & & 1 & & 1 & & 1 & & 1 & & & \\ & & & 1 & & 2 & & 2 & & 2 & & 1 & & \\ & & 1 & & 3 & & 4 & & 4 & & 3 & & 1 & \\ & 1 & & 4 & & 7 & & 8 & & 7 & & 4 & & 1  \\ 1 & & 5 & & 11 & & 15 & & 15 & & 11 & & 5 & & 1 \end{tabular}\]

 The sum of a row of Pascal's triangle is $2^{n-1}$; the sum of two of each of these rows, subtracting away the $2$ ones we included, yields $f(n) = 2^n - 2$. Now, $f(100) = 2^{100} - 2 \equiv 2 \pmod{4}$ and $f(100) = 2^{100} - 2 \equiv 2^{20 \cdot 5} - 2 \equiv -1 \pmod{25}$, and by the \href{https://artofproblemsolving.com/wiki/index.php?title=Chinese\_Remainder\_Theorem}{Chinese Remainder Theorem}, we have $f(100) \equiv 74 \pmod{100} \Longrightarrow \mathrm{(E)}$.

 \textbf{Remark (Chinese Remainder Theorem)} Solution 1 needs to calculate the last $2$ digits of $2^{100}$, meaning to solve $2^{100} (\bmod{ 100})$. $100 = 4 \cdot 25$, we are going to get $n \bmod{ 100}$ from $n \bmod{ 4}$ and $n \bmod{ 25}$.

 By Chinese Remainder Theorem, the general solution of the system of $2$ linear congruences is:

 \begin{verbatim}
$n \equiv r_1 (\bmod { \quad m_1})$, $n \equiv r_2 (\bmod { \quad m_2})$, $(m_1, m_2) = 1$Find $k_1$ and $k_2$ such that $k_1 m_1 \equiv 1 (\bmod{ \quad m_2})$, $k_2 m_2 \equiv 1 (\bmod{ \quad m_1})$Then $n \equiv k_2 m_2 r_1 + k_1 m_1 r_2 (\bmod{ \quad m_1 m_2})$
\end{verbatim}

In this problem, $2^{100} \equiv (2^{10})^{10} \equiv 1024^{10} \equiv 24^{10} \equiv (-1)^{10} \equiv 1 (\bmod{ \quad 25})$, $2^{100} \equiv 0 (\bmod{ \quad 4})$:

 \begin{verbatim}
$n \equiv 1 (\bmod{ \quad 25})$, $n \equiv 0 (\bmod{ \quad 4})$, $(25, 4) = 1$$1 \cdot 25 \equiv 1 (\bmod{ \quad 4})$, $19 \cdot 4 \equiv 1 (\bmod{ \quad 25})$Then $n \equiv 19 \cdot 4 \cdot 1 + 1 \cdot 25 \cdot 0 \equiv 76 (\bmod{ \quad 100})$
\end{verbatim}





\textbf{Solution 2 (induction)}

We sum the first few rows: $0, 2, 6, 14, 30, 62$. They are each two less than a power of $2$, so we try to prove it:

 Let the sum of row $n$ be $S_n$. To generate the next row, we add consecutive numbers. So we double the row, subtract twice the end numbers, then add twice the end numbers and add two. That makes $S_{n+1}=2S_n-2(n-1)+2(n-1)+2=2S_n+2$. If $S_n$ is two less than a power of 2, then it is in the form $2^x-2$. $S_{n+1}=2^{x+1}-4+2=2^{x+1}-2$. 

 Since the first row is two less than a power of 2, all the rest are. Since the sum of the elements of row 1 is $2^1-2$, the sum of the numbers in row $n$ is $2^n-2$. Thus, using \href{https://artofproblemsolving.com/wiki/index.php?title=Modular\_arithmetic}{Modular arithmetic}, $f(100)=2^{100}-2 \bmod{100}$. $2^{10}=1024$, so $2^{100}-2\equiv 24^{10}-2\equiv (2^3 \cdot 3)^{10} - 2$ $\equiv 1024^3 \cdot 81 \cdot 81 \cdot 9 - 2 \equiv 24^3 \cdot 19^2 \cdot 9 - 2$ $\equiv 74\bmod{100} \Rightarrow \mathrm{(E)}$.

 



\textbf{Solution 3 (plain recurrence solving)}

We derive the recurrence $S_{n+1}=2S_n + 2$ as above. Without guessing the form of the solution at this point we can easily solve this recurrence. Note that one can easily get rid of the "$+2$" as follows: Let $S_n=T_n-2$. Then $S_{n+1}=T_{n+1}-2$ and $2S_n+2 = 2(T_n-2)+2 = 2T_n-2$. Therefore $T_{n+1}=2T_n$. This obviously solves to $T_n=2^{n-1} T_1$. As $S_1=0$, we have $T_1=2$. Therefore $T_n=2^n$ and consecutively $S_n=2^n-2$.

 


	\item (\textit{AHSME 1995}) Two parallel chords in a circle have lengths 10 and 14, and the distance between them is 6. The chord parallel to these chords and midway between them is of length $\sqrt {a}$ where $a$ is

                    \begin{figure}[H]
                        \centering
                        \includegraphics[width=0.5\textwidth]{Problem/AHSME_I_1995_28_Problem_Figure_1.png}
                    \end{figure}
                    
$\mathrm{(A) \ 144 } \qquad \mathrm{(B) \ 156 } \qquad \mathrm{(C) \ 168 } \qquad \mathrm{(D) \ 176 } \qquad \mathrm{(E) \ 184 }$

 \textbf{Answer: }E



\textbf{Solution 1}

We let $O$ be the center, $\overline{A_1AA_2}$, $\overline{B_1BB_2}$ represent the chords with length $10, 14$ respectively (as shown below). Connecting the endpoints of the chords with the center, we have several right triangles. However, we do not know whether the two chords are on the same side or different sides of the center of the circle. 

 
                    \begin{figure}[H]
                        \centering
                        \includegraphics[width=0.5\textwidth]{Solution/AHSME_I_1995_28_Solution_Figure_1.png}
                    \end{figure}
                     By the \href{https://artofproblemsolving.com/wiki/index.php?title=Pythagorean\_Theorem}{Pythagorean Theorem} on $\triangle OBB_1$, we get $x^2 + 7^2 = r^2 \Longrightarrow x = \sqrt{r^2 - 49}$, where $x$ is the length of the other leg. Now the length of the leg of $\triangle OAA_1$ is either $6 + x$ or $6 - x$ depending whether or not $\overline{A_1A_2}, \overline{B_1B_2}$ are on the same side of the center of the circle:

 
 \begin{eqnarray*}(6 \pm \sqrt{r^2 - 49})^2 + 5^2 &=& r^2\\ 12 \pm 12\sqrt{r^2 - 49} &=& 0\end{eqnarray*}

 Only the negative works here (thus the two chords are on \textit{opposite} sides of the center), and solving we get $x=1, r = 5\sqrt{2}$. The leg formed in the right triangle with the third chord is $3 - x = 2$, and by the Pythagorean Theorem again

 
 \[(\sqrt{a}/2)^2 + 2^2 = (5\sqrt{2})^2 \Longrightarrow \sqrt{a} = 184 \Rightarrow \mathrm{(E)}.\]




\textbf{Solution 2}

Let $AB$ be the chord of length $10$, $CD$ be the chord of length $14$, $E$ be the point on $CD$ such that $AE = 6$, $F$ be the point on $CD$ such that $BF = 6$. Then $CE = FD = 2$. Extend $BF$ intersecting the circle at $G$ and let $FG = x$. By the \href{https://artofproblemsolving.com/wiki/index.php?title=Power\_of\_a\_Point\_Theorem}{Power of a Point Theorem},
 \[6*x = 12*2\Rightarrow x = 4.\]Since $\angle ABG$ is a \href{https://artofproblemsolving.com/wiki/index.php?title=Right\_angle}{right angle} and $A$ and $G$ are on the circle, the diameter $AG = 10\sqrt {2}$ by Pythagorean Theorem, so the radius is $5\sqrt {2}$.Let $O$ be the center of the circle, $H$ be the midpoint of $CD$. Then $CH = 7$ and $\angle OHC$ is a right angle. Then $OH = 1$, again by Pythagorean Theorem.The chord midway between $AB$ and $CD$ is a distance of 3 away from each, so 2 away from $O$. Using the Pythagorean Theorem one more time, half the length of the chord is equal to $\sqrt {46}$. Then $\sqrt {a} = 2\sqrt {46}\Rightarrow a = 184$. 

 


	\item (\textit{AHSME 1995}) For how many three-element sets of distinct positive integers $\{a,b,c\}$ is it true that $a \times b \times c = 2310$?


 $\mathrm{(A) \ 32 } \qquad \mathrm{(B) \ 36 } \qquad \mathrm{(C) \ 40 } \qquad \mathrm{(D) \ 43 } \qquad \mathrm{(E) \ 45 }$

 \textbf{Answer: }C



\textbf{Solution 1}

$2310 = 2 \cdot 3 \cdot 5 \cdot 7 \cdot 11$. We wish to figure out the number of ways to distribute these prime factors amongst 3 different integers, without over counting triples which are simply permutations of one another.

 We can account for permutations by assuming WLOG that $a$ contains the prime factor 2. Thus, there are $3^4$ ways to position the other 4 prime numbers. Note that, with the exception of when all of the prime factors belong to $a$, we have over counted each case twice, as for when we put certain prime factors into $b$ and the rest into $c$, we count the exact same case when we put those prime factors which were in $b$ into $c$.

 Thus, our total number of cases is $\frac{3^4 - 1}{2} = 40 \Rightarrow \boxed{C}.$




\textbf{Solution 2}

The prime factorization of $2310$ is $2310 = 2 \cdot 3 \cdot 5 \cdot 7 \cdot 11.$ Therefore, we have the equation 
 \[abc = 2310 = 2 \cdot 3 \cdot 5 \cdot 7 \cdot 11,\]where $a, b, c$ must be distinct positive integers and order does not matter. There are $3$ ways to assign each prime number on the right-hand side to one of the variables $a, b, c,$ which gives $3^5 = 243$ solutions for $(a, b, c).$ However, three of these solutions have two $1$s and one $2310,$ which contradicts the fact that $a, b, c$ must be distinct. Because each prime factor appears only once, all other solutions have $a, b, c$ distinct. Correcting for this, we get $243 - 3 = 240$ ordered triples $(a, b, c)$ where $a, b, c$ are all distinct.

 Finally, since order does not matter, we must divide by $3!,$ the number of ways to order $a, b, c.$ This gives the final answer, 
 \[\frac{240}{3!} = \frac{240}{6} = \boxed{40}.\]



	\item (\textit{AHSME 1995}) A large cube is formed by stacking 27 unit cubes. A plane is perpendicular to one of the internal diagonals of the large cube and bisects that diagonal. The number of unit cubes that the plane intersects is

                    \begin{figure}[H]
                        \centering
                        \includegraphics[width=0.5\textwidth]{Problem/AHSME_I_1995_30_Problem_Figure_1.png}
                    \end{figure}
                    
$\mathrm{(A) \ 16 } \qquad \mathrm{(B) \ 17 } \qquad \mathrm{(C) \ 18 } \qquad \mathrm{(D) \ 19 } \qquad \mathrm{(E) \ 20 }$

 \textbf{Answer: }D



\textbf{Solution 1}

Place one corner of the cube at the origin of the coordinate system so that its sides are parallel to the axes.

 Now consider the diagonal from $(0,0,0)$ to $(3,3,3)$. The midpoint of this diagonal is at $\left(\frac 32,\frac 32,\frac 32\right)$. The plane that passes through this point and is orthogonal to the diagonal has the equation $x+y+z=\frac 92$.

 The unit cube with opposite corners at $(x,y,z)$ and $(x+1,y+1,z+1)$ is intersected by this plane if and only if $x+y+z < \frac 92 < (x+1)+(y+1)+(z+1)=(x+y+z)+3$. Therefore the cube is intersected by this plane if and only if $x+y+z\in\{2,3,4\}$.

 There are six cubes such that $x+y+z=2$: permutations of $(1,1,0)$ and $(2,0,0)$. 
Symmetrically, there are six cubes such that $x+y+z=4$. 
Finally, there are seven cubes such that $x+y+z=3$: permutations of $(2,1,0)$ and the central cube $(1,1,1)$.

 That gives a total of $\boxed{19}$ intersected cubes.

 Note that there are only 8 cubes that are not intersected by our plane: 4 in each of the two opposite corners that were connected by the original diagonal.

 



\textbf{Solution 2}

Place the cube so that its space diagonal is perpendicular to the ground. The space diagonal has length of $3\sqrt{3}$, the altitude of the top vertex of the newly placed cube is $3\sqrt{3}$. The plane perpendicular and bisecting the space diagonal is now parallel to the ground and also bisecting the space diagonal into $\frac{3\sqrt{3}}{2}$, so that the height of the plane is $\frac{3\sqrt{3}}{2}$.

 By symmetry, the space diagonal is trisected by the pyramid at the top of the cube and the pyramid at the bottom of the cube.

 We can prove that the space diagonal is trisected. Let the altitude of the pyramid at the top of the cube be $h$. The base of the pyramid is an equilateral triangle with side length of $\sqrt{3^2+3^2}=3\sqrt{2}$. The height of the triangle is $\frac{ \sqrt{3} }{2} \cdot 3\sqrt{2}$. The distance of the center of the triangle to the vertex is $\frac23 \cdot \frac{ \sqrt{3} }{2} \cdot 3\sqrt{2} = \sqrt {6}$. Therefore, $h = \sqrt{3^2-(\sqrt {6})^2} = \sqrt{9-6}=\sqrt{3}$. The altitude of the pyramid at the bottom of the cube is also $h$. The altitude in the middle is $3\sqrt{3}-\sqrt{3}-\sqrt{3}=\sqrt{3}$.

 The altitude of the vertex at the top is $3\sqrt{3}$. The altitude of the second highest $3$ vertices are all $2\sqrt{3}$. The altitude of the third highest $3$ vertices are all $\sqrt{3}$. The altitude of the bottom-most vertex is $0$.

 By scale, for the unit cube, place the cube so that its space diagonal is perpendicular to the ground. The altitude of the vertex at the top is $\sqrt{3}$. The altitude of the second highest $3$ vertices are all $\frac{2\sqrt{3}}{3}$. The altitude of the third highest $3$ vertices are all $\frac{\sqrt{3}}{3}$. The altitude of the bottom-most vertex is $0$.

 The length of the space diagonal of a unit cube is $\sqrt{3}$. The highest vertex of the bottom-most unit cube has an altitude of $\sqrt{3}$. As $\sqrt{3} < \frac{3\sqrt{3}}{2}$, therefore, the plane will not pass through the unit cube at the bottom.

 For the next $3$ cubes from the bottom, the altitude of their highest vertex is $\frac{\sqrt{3}}{3} + \sqrt{3} = \frac{4\sqrt{3}}{3}$. As $\frac{4\sqrt{3}}{3} < \frac{3\sqrt{3}}{2}$, therefore, the plane will not pass through the next $3$ unit cubes.

 For the next $3$ cubes from the bottom, the altitude of their highest vertex is $\frac{2\sqrt{3}}{3} + \frac{\sqrt{3}}{3} = \frac{5\sqrt{3}}{3}$. As $\frac{5\sqrt{3}}{3} > \frac{3\sqrt{3}}{2}$, therefore, the plane will pass through the next $3$ unit cubes.

 So at the bottom half of the cube, there are only $1+3 = 4$ unit cubes that the plane does not passes through. By symmetry, the plane will not pass through $4$ unit cubes at the top half of the cube.

 Thus, the plane does not pass through $4+4 = 8$ unit cubes, it passes through $27-8=\boxed{\textbf{(D) } 19}$ unit cubes.

 


\end{enumerate}